% Part: first-order-logic
% Chapter: introduction
% Section: semantic-notions

\documentclass[../../../include/open-logic-section]{subfiles}

\begin{document}

\olfileid{fol}{int}{sem}

\olsection{અર્થવિચારી ખ્યાલો}

ઉપર આપણે કહ્યું કે $\Sat{M}{!A}[s]$ છે કે નહિ એ વિચારતી વખતે સગવડ
માટે $s$ બધાં !!{variable}sને મૂલ્યો નિયુક્ત કરે એવું રાખીએ છીએ, પરંતુ
વપરાય છે માત્ર~$!A$માંના !!{variable}sને તે નિયુક્ત કરેલાં મૂલ્યો.
હકીકતમાં,~$!A$માંના માત્ર \emph{મુક્ત} ચલોને આપેલાં મૂલ્યો જ મહત્ત્વ
ધરાવે છે. આપણે ચોકસાઈ રાખીએ છીએ, તેથી આ હકીકત અલબત્ત સાબિત કરીશું.
!!{sentence}sમાં કોઈ મુક્ત ચલ નથી; તેથી તેઓ !!a{structure}માં સંતોષાય
છે કે નહિ તેમાં~$s$નો જરાય ફેર પડતો નથી. એટલે $!A$ !!a{sentence} હોય
ત્યારે આપણે $\Sat{M}{!A}$નો અર્થ ``બધાં~$s$ માટે
$\Sat{M}{!A}[s]$'' એમ વ્યાખ્યાયિત કરી શકીએ છીએ; અને હકીકતમાં આ ત્યારે
અને માત્ર ત્યારે સત્ય છે જ્યારે ઓછામાં ઓછી એક~$s$ માટે
$\Sat{M}{!A}[s]$ હોય. !!{formula}s માટે સંતોષની કામ કરતી વ્યાખ્યા
આપવા !!{variable}-નિયુક્તિઓ દાખલ કરવી પડે છે, પરંતુ !!{sentence}s
માટેનો સંતોષ !!{variable}-નિયુક્તિઓથી સ્વતંત્ર છે.

``$\Sat{M}{!A}$''ની વ્યાખ્યા મળી જાય પછી પ્રથમ-ક્રમ તર્કશાસ્ત્રનાં
!!{sentence}s માટે ``કિસ્સો'' અને ``એમાં સત્ય''નો અર્થ જાણી લઈએ છીએ.
!!{sentence}s માટે $\Sat{M}{!A}$ની વ્યાખ્યાને આધારે પછી પ્રામાણ્ય,
ફલિતતા અને સંતોષ્યતાના મૂળ અર્થવિચારી ખ્યાલો વ્યાખ્યાયિત કરી શકીએ
છીએ. દરેક !!{structure} કોઈ વાક્યને સંતોષે, એટલે $\Entails !A$ હોય,
તો તે વાક્ય પ્રમાણભૂત છે. !!{sentence}sના કોઈ ગણથી વાક્ય ફલિત થાય,
એટલે $\Gamma \Entails !A$ હોય, ત્યારે અને માત્ર ત્યારે જ્યારે
~$\Gamma$નાં બધાં !!{sentence}sને સંતોષતી દરેક !!{structure},~$!A$ને
પણ સંતોષે. અને કોઈ એક !!{structure}, !!{sentence}sના ગણમાંનાં બધાં
!!{sentence}sને એકસાથે સંતોષે, તો એ ગણ સંતોષ્ય છે.

!!{formula}s અનુમાનાત્મક રીતે વ્યાખ્યાયિત છે અને સંતોષ પણ
!!{formula}sની રચના પર અનુમાન વડે વ્યાખ્યાયિત થાય છે; તેથી આપણા
અર્થવિચારના ગુણધર્મો સાબિત કરવા અને વ્યાખ્યાયિત અર્થવિચારી ખ્યાલોને
પરસ્પર જોડવા આપણે અનુમાન વાપરી શકીએ છીએ. એમાંથી કેટલાક ગુણધર્મો
એકત્ર કરીને સાબિત કરીશું; અંશતઃ તેઓ પોતપોતાની રીતે રસપ્રદ છે, પરંતુ
મુખ્યત્વે એટલા માટે કે આગળ અર્થવિચાર અને !!{derivation} તંત્રો વચ્ચેનો
સંબંધ તપાસતી વખતે તેમાંના ઘણા ઉપયોગી થશે. એ કરવા માટે હજુ ઉલ્લેખ
નથી કર્યો એવા કેટલાક વાક્યરચનાત્મક ખ્યાલો અને ક્રિયાઓ પણ (ચોક્કસ
રીતે, એટલે કે અનુમાન વડે) વ્યાખ્યાયિત કરવા પડશે.

\end{document}
